% ID: SEL-D0010-BD-E08
% ESTADO: cotejada. Segunda pasada visual completa y borrador compilado revisados.
% Fuente: Selectividad/matematicas_aplicadas/EXTRA_TITULAR_MAT-APLIC-CCSS-JULIO-IMPRESO.pdf
% SHA-256: 548f60ee334516c7fa918dba3c28fb3ea5fe753cf76eb05402d7b8acb5a3b560
% Página: 2; bloque D; ejercicio original 8.

Sea $X$ una variable aleatoria que sigue una ley Normal de media poblacional
desconocida y desviación típica 4.
\begin{enumerate}[label=\alph*)]
  \item \textbf{(0.5 puntos)} \textquestiondown{}Cuál es la desviación típica
  de la distribución de las medias de las muestras de tamaño 12 de la variable
  aleatoria $X$?
  \item \textbf{(1 punto)} Para estimar la media poblacional de la variable
  $X$, se toma una muestra aleatoria de tamaño 12, obteniéndose los siguientes
  resultados:
  \[
  11.8\quad 10\quad 9.8\quad 12\quad 9.7\quad 10.8\quad
  9.6\quad 11.3\quad 10.4\quad 12.2\quad 9.1\quad 10.5
  \]
  Con los datos obtenidos de la muestra, determine un intervalo de confianza
  al 97\% para estimar la media poblacional.
  \item \textbf{(1 punto)} Calcule el tamaño mínimo que debe tener una muestra,
  para que, con el mismo nivel de confianza, el error cometido al estimar la
  media poblacional sea menor que 1.2.
\end{enumerate}
